\subsection{逼近空间}

给定局部凸空间 E 和 F，回顾（EVT, III, p. 14, 例 2），我们用 \(\mathcal{L}_{\mathrm{pc}}(\mathrm{E};\mathrm{F})\) 表示给 \(\mathcal{L}(\mathrm{E};\mathrm{F})\) 赋以预紧收敛拓扑所得到的局部凸空间。我们也用 \(\mathcal{L}_{\mathrm{pc}}(\mathrm{E})\) 表示局部凸空间 \(\mathcal{L}_{\mathrm{pc}}(\mathrm{E};\mathrm{E})\)。当局部凸空间 E 是豪斯多夫的且拟完备（EVT, III, p. 8, def. 6）时，\(\mathcal{L}(\mathrm{E};\mathrm{F})\) 上的预紧收敛拓扑与紧收敛拓扑重合，因为 E 的预紧子集的闭包是紧的（EVT, III, p. 8）。

定义 2. --- 局部凸空间 E 称为具有逼近性质，或又称为逼近空间，如果 E 的恒等映射 \(1_{E}\) 在空间 \(\mathcal{L}_{\mathrm{pc}}(\mathrm{E})\) 中附着于 E 的连续有限秩自同态所成的集 \(\mathcal{L}^{\mathrm{f}}(\mathrm{E})\)。

特别地，每个有限维局部凸空间都是逼近空间。

注记. --- 1) 为使局部凸空间 E 和 F 的拓扑直和是逼近空间，必须且只需 E 和 F 都是逼近空间。

\begin{enumerate}
\def\labelenumi{\arabic{enumi})}
\setcounter{enumi}{1}
\item
  设 E 是局部凸空间。设 N 是 \(\{0\}\) 在 E 中的闭包，P 是 N 在 E 中的一个代数补子空间。那么 P 是 N 在 E 中的拓扑补，且同构于局部凸空间 E/N。空间 N 是逼近空间。由注记 1，为使 E 是逼近空间，必须且只需豪斯多夫局部凸空间 E/N 是逼近空间。
\item
  设 E 是局部凸空间，A 是 \(\mathcal{L}^{\mathrm{f}}(\mathrm{E})\) 的一个等度连续子集，T 是 E 的一个完全子集。若 \(1_{\mathrm{E}}\) 对 T 上的逐点收敛拓扑附着于 A，则 \(1_{\mathrm{E}}\) 在 \(\mathcal{L}_{\mathrm{pc}}(\mathrm{E})\) 中附着于 A（EVT, III, p. 16, prop. 4 及 p. 17, prop. 5），从而 E 是逼近空间。
\item
  设 E 是 C 上的局部凸空间。用 \(E_{0}\) 表示 E 的底层 R 上的局部凸空间。为使 E 是逼近空间，必须且只需 \(E_{0}\) 本身是逼近空间。事实上，该条件是必要的；我们来证明它是充分的。于是设 \(E_{0}\) 是逼近空间。设 C 是 E 的一个预紧子集，U 是 E 中 0 的一个凸均衡邻域。令 \(C' = C \cup iC\)。存在一个从 \(E_{0}\) 到 \(E_{0}\) 的连续有限秩 R-线性映射 u，使得 \(x - u(x)\) 对一切 \(x \in C'\) 都属于 U。对 E 中的一切 x，令 \(v(x) = \frac{1}{2}(u(x) - iu(ix))\)。这样就定义了一个从 E 到 E 的连续有限秩 C-线性映射。对每个 \(x \in C\)，我们有 \(x \in C'\)，\(ix \in C'\)，并且 \(x - v(x) = \frac{1}{2}(x - u(x)) - \frac{i}{2}(ix - u(ix))\)，从而 \(x - v(x)\) 属于 U。这就证明了 E 是逼近空间。
\item
  设 E 是 R 上的局部凸空间。为使复化局部凸空间 \(\mathrm{E}_{(\mathbf{C})}\) 是逼近空间，必须且只需 E 本身是逼近空间。这由注记 1 和注记 4 得出，因为 \(\mathrm{E}_{(\mathbf{C})}\) 的底层实局部凸空间同构于 \(E \times E\)。
\end{enumerate}

命题 10. --- 设 E 是希尔伯特空间。那么 E 是逼近空间。

E 上的有限秩正交投影算子所成的集 A 是等度连续的。设 \(n \geqslant 1\) 是整数，\((x_{1}, \ldots, x_{n})\) 是 E 的元素的一个族。用 V 表示由诸 \(x_{i}\) 生成的向量子空间，用 \(p_{V}\) 表示以 V 为像的正交投影算子。我们有 \(p_{\mathrm{V}}(x_{i}) = x_{i}\)（对 \(1 \leqslant i \leqslant n\)）。这就蕴含 \(1_{E}\) 对逐点收敛拓扑附着于 A，从而 E 是逼近空间（注记 3）。

逼近空间的其他例子将在 III, p. 20, 第 7 目中给出。

引理 2. --- 设 E 是局部凸空间。那么 E 是逼近空间，当且仅当：对 E 的每个预紧子集 C 和 E 中 0 的每个邻域 U，存在一个整数 \(n \geqslant 0\)、E 的元素 \(e_{1}, \ldots, e_{n}\) 以及 E 上的连续线性形式 \(f_{1}, \ldots, f_{n}\)，使得 \(x - \sum_{i=1}^{n} f_{i}(x)e_{i}\) 对一切 \(x \in C\) 都属于 U。

预紧收敛拓扑的定义表明该条件是充分的。反之，设 E 是逼近空间。设 P 是 \(\{0\}\) 的闭包在 E 中的一个拓扑补空间（参见注记 2）。此时，\(u \in \mathcal{L}^{\mathrm{f}}(\mathrm{E})\)（满足 \(u(\mathrm{E}) \subset \mathrm{P}\)）所成的集 A 在 \(\mathcal{L}^{\mathrm{f}}(\mathrm{E})\) 中稠密。由于 P 是分离的，A 的每个元素 \(u\) 都形如 \(x \mapsto \sum_{i=1}^{n} f_i(x)e_i\)，其中 \(n \in \mathbf{N}\)，\(e_1, \ldots, e_n\) 是 E 的元素，而 \(f_1, \ldots, f_n\) 是 E 上的连续线性形式（EVT, I, p. 14, th. 2）。这就证明了该条件是必要的。

注记 6. --- 存在不具有逼近性质的巴拿赫空间，这是由 P. ENFLO 证明的（A counterexample to the approximation problem in Banach spaces, Acta Math. 130 (1973), p. 309--317; 参见 III, p. 112, 习题 25）。这回答了 S. BANACH 的一个问题（Théorie des opérations linéaires, Monografie Matematyczne I, Warszawa, 1932, 第 VI 章 § 1 中的注记）。另见 p. 21 上的注记。

命题 11. --- 设 E 和 F 是局部凸空间。假设 E 是逼近空间。

\begin{enumerate}
\def\labelenumi{\alph{enumi})}
\item
  集 \(\mathcal{L}^{\mathrm{f}}(\mathrm{E};\mathrm{F})\) 在 \(\mathcal{L}_{\mathrm{pc}}(\mathrm{E};\mathrm{F})\) 中稠密；
\item
  集 \(\mathcal{L}^{\mathrm{f}}(\mathrm{F};\mathrm{E})\) 在 \(\mathcal{L}_{\mathrm{pc}}(\mathrm{F};\mathrm{E})\) 中稠密；
\item
  设 \(\mathfrak{S}\) 是 F 的有界子集所成的一个集，\(u \in \mathcal{L}(\mathrm{F};\mathrm{E})\)。假设在 \(u\) 下，属于 \(\mathfrak{S}\) 的 F 的每个子集的像都是预紧的。那么 \(u\) 附着于 \(\mathcal{L}^{\mathrm{f}}(\mathrm{F};\mathrm{E})\)（对 \(\mathfrak{S}\)-拓扑）（EVT, III, p. 13）。
\end{enumerate}

设 \(v\) 是 \(\mathcal{L}(\mathrm{E};\mathrm{F})\) 的一个元素。映射 \(\varphi\colon w\mapsto v\circ w\)（从 \(\mathcal{L}_{\mathrm{pc}}(\mathrm{E})\) 到 \(\mathcal{L}_{\mathrm{pc}}(\mathrm{E};\mathrm{F})\) 中）是连续的（TG, X, p. 5, prop. 3）。我们有 \(\varphi(1_{\mathrm{E}})=v\) 及 \(\varphi(\mathcal{L}^{\mathrm{f}}(\mathrm{E}))\subset\mathcal{L}^{\mathrm{f}}(\mathrm{E};\mathrm{F})\)，因此 \(v\) 附着于 \(\mathcal{L}^{\mathrm{f}}(\mathrm{E};\mathrm{F})\)（在 \(\mathcal{L}_{\mathrm{pc}}(\mathrm{E};\mathrm{F})\) 中）。这就证明了 \(a\)）。

同样，在 \(c\)）的假设之下，映射 \(\psi: w \mapsto w \circ u\)（从 \(\mathcal{L}_{\mathrm{pc}}(\mathrm{E})\) 到 \(\mathcal{L}_{\mathfrak{S}}(\mathrm{F};\mathrm{E})\) 中）是连续的（同前引处）。我们有 \(\psi(1_{\mathrm{E}}) = u\) 及 \(\psi(\mathcal{L}^{\mathrm{f}}(\mathrm{E})) \subset \mathcal{L}^{\mathrm{f}}(\mathrm{F};\mathrm{E})\)，因此 \(u\) 附着于 \(\mathcal{L}^{\mathrm{f}}(\mathrm{F};\mathrm{E})\)（在 \(\mathcal{L}_{\mathfrak{S}}(\mathrm{F};\mathrm{E})\) 中）。

F 的预紧子集在从 F 到 E 的连续线性映射下的像是预紧的，故 b) 由 c) 得出。

推论. --- 设 E 是分离的逼近空间，F 是局部凸空间。从 F 到 E 的每个紧线性映射都附着于 \(\mathcal{L}^{\mathrm{f}}(\mathrm{F};\mathrm{E})\)（在赋有有界收敛拓扑的空间 \(\mathcal{L}(\mathrm{F};\mathrm{E})\) 中）。

这由命题 11 的 c) 得出，因为 F 的有界子集在从 F 到 E 的紧线性映射下的像在 E 中相对紧（III, p. 2, 注记 1），从而预紧。

命题 12. --- 设 E 是局部凸空间，I 是一个集合，且对每个 \(i \in I\)，设 \(F_{i}\) 是局部凸空间，\(v_{i}: E \to F_{i}\) 是有稠密像的连续线性映射。假设对 E 中 0 的每个邻域 U，存在 \(i \in I\) 及 \(F_{i}\) 中 0 的一个邻域 V，使得 \(v_{i}^{-1}(V) \subset U\)。若诸 \(F_{i}\) 都是逼近空间，则 E 也是逼近空间。

设 A 是 E 的一个预紧子集，U 是 E 中 0 的一个邻域。按假设，存在 \(i \in I\) 及一个连续半范数 \(p\)（定义在 \(F_i\) 上），使得 U 包含 \((p \circ v_i)^{-1}([0,1])\)。令 \(F = F_i\)，\(v = v_i\)，\(B = v(A)\)，并设 F 是逼近空间。集合 B 在 F 中预紧。因此，由引理 2，存在一个整数 \(n \geqslant 1\)、F 的元素 \(y_1, \ldots, y_n\) 以及 F 上的连续线性形式 \(f_1, \ldots, f_n\)，使得对一切 \(y \in B\)，有

\[
p \left(y - \sum_ {j = 1} ^ {n} f _ {j} (y) y _ {j}\right) \leqslant \frac {1}{2}.
\]

由于 B 有界（EVT, III, p. 3, prop. 2），存在实数 M > 0，使得 \(|f_{j}(y)| \leqslant M\) 对一切满足 \(1 \leqslant j \leqslant n\) 的 j 及每个 \(y \in B\) 都成立。此外，由于 \(v(\mathrm{E})\) 在 F 中稠密，对每个满足 \(1 \leqslant j \leqslant n\) 的整数 j，存在 E 的一个元素 \(x_{j}\)，使得 \(p(y_{j} - v(x_{j})) \leqslant (2n\mathrm{M})^{-1}\)。线性映射

\[
u \colon x \longmapsto \sum_ {j = 1} ^ {n} f _ {j} (v (x)) x _ {j}
\]

属于 \(\mathcal{L}^{\mathrm{f}}(\mathrm{E})\)。对每个 \(x\in \mathrm{A}\)，我们有

\[
v (x - u (x)) = v (x) - \sum_ {j = 1} ^ {n} f _ {j} (v (x)) y _ {j} + \sum_ {j = 1} ^ {n} f _ {j} (v (x)) \left(y _ {j} - v \left(x _ {j}\right)\right),
\]

由此 \(p(v(x - u(x))) \leqslant \frac{1}{2} + \frac{nM}{2nM} = 1\)，从而 \(x - u(x) \in U\)。命题得证。

推论 1. --- 逼近空间的一个稠密向量子空间是逼近空间。

推论 2. --- 若局部凸空间 E 的分离完备化是逼近空间，则 E 是逼近空间。

推论 3. --- 一族逼近空间的乘积是逼近空间。

事实上，设 \((\mathrm{E}_{i})_{i\in\mathrm{I}}\) 是一族逼近空间。对 I 的每个有限子集 J，令 \(E_{J}=\prod_{i\in J}E_{i}\)，并用 \(v_{J}\) 表示从 \(E=\prod_{i\in I}E_{i}\) 到 \(E_{J}\) 中的典范映射。局部凸空间 \(E_{J}\) 是逼近空间（注记 1）。于是，把命题 12 应用于局部凸空间 E、局部凸空间族 \((\mathrm{E}_{\mathrm{J}})\) 以及连续线性映射 \(v_{J}:E\to E_{J}\)，即得该推论。

推论 4. --- 设 E 是局部凸空间。若 E 上的每个连续半范数都被一个连续预希尔伯特半范数所控制，则 E 是逼近空间。

设 \(\mathcal{P}\) 是 E 上连续预希尔伯特半范数所成的集。对每个 \(p \in \mathcal{P}\)，假设蕴含：半赋范空间 \(\mathrm{E}_p\)（给 E 赋以半范数 \(p\) 而得到）是逼近空间（推论 2 与命题 10），并且从 E 到 \(\mathrm{E}_p\) 的恒等映射是连续的。因此该推论由命题 12 得出。

引理 3. --- 设 E 是可度量化的局部凸空间，\((x_{n})_{n\in\mathbf{N}}\) 是 E 中收敛于 0 的元素的一个序列。则存在 E 中收敛于 0 的元素的一个序列 \((y_{n})_{n\in\mathbf{N}}\)，以及区间 [0,1] 中收敛于 0 的元素的一个序列 \((\lambda_{n})_{n\in\mathbf{N}}\)，使得 \(x_{n}=\lambda_{n}y_{n}\) 对一切 \(n\in N\) 成立。

由于 E 可度量化，存在 E 中 0 的均衡邻域的一个基本系 \((\mathrm{V}_{m})_{m\in\mathbf{N}}\)，使得 \(V_{0}=E\)，且 \(2^{m+1}V_{m+1}\subset V_{m}\) 对一切 \(m\geqslant0\) 成立。由于 \((x_{n})\) 收敛于 0，存在整数的一个严格递增序列 \((\mathrm{N}_{m})_{m\in\mathbf{N}}\)，使得 \(N_{0}=0\)，且对一切 \(m\geqslant0\)，有 \(x_{n}\in V_{m}\)（对 \(n\geqslant N_{m}\)）。对每个整数 \(n\geqslant0\)，存在唯一的整数 \(m\geqslant0\)，使得 \(N_{m}\leqslant n<N_{m+1}\)，这时我们令 \(\lambda_{n}=2^{-m}\)，\(y_{n}=2^{m}x_{n}\)。如此定义的序列 \((\lambda_{n})\) 收敛于 0。此外，由于 \(y_{n}\in V_{m}\)（对 \(n\geqslant N_{m+1}\)），序列 \((y_{n})\) 在 E 中收敛于 0。最后，对一切 n，我们有 \(x_{n}=\lambda_{n}y_{n}\)。

回顾（EVT, II, p. 28），若 E 是向量空间而 L 是 E 的一个凸均衡子集，我们用 \(E_{L}\) 表示由 L 生成的 E 的向量子空间，赋以以 L 为单位球的半范数。当集合 L 不含任何直线时，这个半范数是一个范数。

引理 4. — 设 E 是弗雷歇空间，A 是 E 的一个紧子集。则存在 E 的一个包含 A 的紧凸均衡子集 L，使得由 E 与由 \(E_{L}\) 在 A 上诱导的拓扑重合。

集合 A 包含于 E 中收敛于 0 的元素的一个序列 \((x_{n})_{n\in\mathbf{N}}\) 的诸点的闭凸均衡包中（EVT, IV, p. 24, cor. 1）。设 \((y_{n})_{n\in\mathbf{N}}\) 与 \((\lambda_{n})_{n\in\mathbf{N}}\) 是满足引理 3 的结论的序列。序列 \((y_{n})\) 的诸点的闭凸均衡包 L 包含 A，且是 E 的一个紧子集（EVT, III, p. 8）。因此，\(E_{L}\) 是巴拿赫空间（EVT, III, p. 8, cor.）。在这个巴拿赫空间中，\(x_{n}\) 的范数被 \(\lambda_{n}\) 所控制，故序列 \((x_{n})\) 趋于 0。\(\widetilde{A}\)（序列 \((x_{n})\) 在 \(E_{L}\) 中的闭凸均衡包）是 \(E_{L}\) 的一个紧子集（EVT, III, p. 8）。由于 L 是 E 的有界子集，\(E_{L}\) 到 E 中的典范单射是连续的。在 \(\widetilde{A}\) 上，由 \(E_{L}\) 诱导的拓扑与由 E 诱导的拓扑重合，且 \(\widetilde{A}\) 是 E 的一个紧子集，从而是闭子集。由于集合 \(\widetilde{A}\) 是凸的、均衡的且包含序列 \((x_{n})\)，故它包含 A。证明完毕。

定理 4. --- 设 E 是弗雷歇空间。

\begin{enumerate}
\def\labelenumi{\alph{enumi})}
\item
  假设 E 是逼近空间。那么，对每个半赋范空间 F，空间 \(\mathcal{L}^{\mathrm{f}}(\mathrm{F};\mathrm{E})\) 对有界收敛拓扑在 \(\mathcal{L}^{\mathrm{c}}(\mathrm{F};\mathrm{E})\) 中稠密；
\item
  反之，假设对每个巴拿赫空间 F，从 F 到 E 的每个紧线性映射都对有界收敛拓扑附着于 \(\mathcal{L}^{\mathrm{f}}(\mathrm{F};\mathrm{E})\)。那么 E 是逼近空间。
\end{enumerate}

设 F 是半赋范空间。\(\mathcal{L}^{\mathrm{f}}(\mathrm{F};\mathrm{E})\) 在 \(\mathcal{L}(\mathrm{F};\mathrm{E})\) 中的闭包含于 \(\mathcal{L}^{\mathrm{c}}(\mathrm{F};\mathrm{E})\)（III, p. 4, 命题 1 与命题 2）。若 E 是逼近空间，则由命题 11 的推论，它等于 \(\mathcal{L}^{\mathrm{c}}(\mathrm{F};\mathrm{E})\)。这就证明了断言 \(a\)）。

假设 \(b\)）的假设成立。设 \(\varepsilon > 0\) 是实数。设 A 是 E 的一个紧子集，\(p\) 是 E 上的一个连续半范数。设 L 是 E 的一个紧凸均衡子集，使得 A 是赋范空间 \(\mathrm{E_L}\) 的一个紧子集（引理 4）。典范单射 \(j: \mathrm{E_L} \to \mathrm{E}\) 是紧的，且 \(\mathrm{E_L}\) 是巴拿赫空间（EVT, III, p. 8, cor.）。因此，按假设，存在一个整数 \(n \geqslant 1\)、E 的元素 \(e_1, \ldots, e_n\) 以及连续线性形式 \(\ell_1, \ldots, \ell_n\)（定义在 \(\mathrm{E_L}\) 上），使得映射 \(v\)（它从 \(\mathrm{E_L}\) 到 E，由 \(v(x) = \sum_{i=1}^{n} \ell_i(x)e_i\) 定义）满足 \(p(x - v(x)) \leqslant \frac{\varepsilon}{2}\)（对 \(x \in A\)）。映射 \({}^{t}j: E' \to (E_{L})'\) 的像对弱拓扑在 \((E_{L})'\) 中稠密（EVT, IV, p. 6, prop. 5）。在 \((E_{L})'\) 上，紧收敛拓扑与 \((E_{L})'\) 和 \(E_{L}\) 之间的对偶相容（EVT, IV, p. 3, 例）。因此 \({}^{t}j(E')\) 对紧收敛拓扑在 \((E_{L})'\) 中稠密（EVT, IV, p. 1, prop. 1），从而存在 E 上的连续线性形式 \(f_{1}, \ldots, f_{n}\)，使得

\[
| \ell_ {i} (x) - f _ {i} (x) | p (e _ {i}) \leqslant \frac {\varepsilon}{2 n}
\]

其中 \(x \in \mathrm{A}\)，\(1 \leqslant i \leqslant n\)。自同态 \(u \colon x \mapsto \sum_{i=1}^{n} f_i(x)e_i\)（\(\mathrm{E}\) 的）属于 \(\mathcal{L}^{\mathrm{f}}(\mathrm{E})\)，且对每个 \(x \in \mathrm{A}\)，我们有

\[
p (x - u (x)) \leqslant p (x - v (x)) + p (v (x) - u (x)) \leqslant \frac {\varepsilon}{2} + n \times \frac {\varepsilon}{2 n} = \varepsilon .
\]

由此推出，\(1_{E}\) 对紧收敛拓扑附着于 \(\mathcal{L}^{\mathrm{f}}(\mathrm{E})\)。而由于 E 完备，这个拓扑与预紧收敛拓扑重合。因此 E 是逼近空间。

